% Aula 01 --- risco condicional x risco esperado, e a barra de erro do teste.
A palavra "risco" apareceu durante o nosso curso com dois sentidos, e a diferença entre eles
não é preciosismo: ela decide qual pergunta o seu número está respondendo.

Quando $g$ é tratada como uma função \emph{fixa}, a esperança em
$\E[(Y-g(\X))^2]$ percorre só a observação nova $(\X,Y)$. Chamamos isso de
\textbf{risco condicional}. Quando $g=\ghat$ é tratada como aleatória --- e ela é,
porque foi construída a partir de $\Dados$ ---, a esperança percorre também as
amostras de treino possíveis, e temos o \textbf{risco esperado}.
\begin{itens}
  \item \pts{1,3} Cada um dos dois oferece uma garantia diferente. Descreva as
        duas, e diga qual interessa a quem já ajustou um modelo e vai colocá-lo em
        produção, e qual interessa a quem ainda está escolhendo entre diversos métodos.
  \item \pts{1,2} Você separou um conjunto de teste de $m$ observações que não
        participou de nenhum treinamento, e calculou $\riscohat(\ghat)=\frac1m\sum_{k=1}^m W_k$,
        com $W_k=(Y_k-\ghat(\X_k))^2$. Qual dos dois riscos esse número estima?
        Para responder, pergunte-se o que está fixo e o que varia no momento em
        que você faz a conta.
  \item \pts{1,5} Um conjunto de teste permite algo que a validação cruzada não
        permite: uma barra de erro honesta. Explique por que os $W_k$ são i.i.d.\ e
        escreva o intervalo de confiança aproximado de 95\% que o Teorema Central
        do Limite fornece.
\end{itens}

\begin{solucao}
\textbf{(a)} O \textbf{risco condicional} é uma garantia sobre \emph{aquela
função específica}: se eu usar este $\ghat$, já ajustado, em observações novas, meu
erro quadrático médio esperado é tanto.

O \textbf{risco esperado} é uma garantia sobre o \emph{procedimento}: se eu
sortear uma amostra de treino, rodar este método e usar o que sair, meu erro
esperado é tanto. A esperança percorre também os treinos possíveis.

Quem vai pôr o modelo em produção quer o \textbf{condicional}. O modelo dele é um
só, já existe, não vai ser reajustado --- saber como se comportam os modelos que
aquele método \emph{costuma} produzir não paga as contas.

Quem está comparando métodos quer o \textbf{esperado}. A pergunta ali é "qual
deles tende a produzir bons modelos", e uma rodada de sorte com um método ruim não
deveria decidir a escolha.

\smallskip
\textbf{(b)} Estima o \textbf{risco condicional}.

No momento da conta, $\ghat$ já foi ajustado e não muda: é uma função fixa. O que
varia de uma parcela para outra é apenas o par $(\X_k,Y_k)$ sorteado para o
teste. A aleatoriedade do treino --- que outra amostra teria dado outro $\ghat$
--- não entra na conta em momento nenhum, porque o treino aconteceu antes e não é
repetido.

Dito de outro modo: para estimar o risco \emph{esperado} você teria de repetir o
processo inteiro, sorteando vários treinos, e não é o que um único conjunto de
teste faz.

\smallskip
\textbf{(c)} Condicionalmente a $\ghat$ --- que é a situação, por (b) --- cada
$W_k=(Y_k-\ghat(\X_k))^2$ é função de um único par $(\X_k,Y_k)$. Esses pares são
i.i.d.\ entre si e independentes do treino, e a mesma função $\ghat$ é aplicada a
todos. Logo os $W_k$ são i.i.d., com média $\risco(\ghat)$.

Como $\riscohat$ é média de i.i.d., o Teorema Central do Limite se aplica e dá
\[
  \riscohat(\ghat)\;\pm\;2\sqrt{\widehat\sigma^2/m},
  \qquad
  \widehat\sigma^2=\frac1m\sum_{k=1}^m(W_k-\overline W)^2 .
\]

Há uma hipótese fazendo todo o trabalho aqui: o teste não participou de nada. Se ele
tivesse sido usado para escolher qualquer coisa, os $W_k$ deixariam de ser
independentes do $\ghat$ e o intervalo perderia o sentido.
\end{solucao}
